% GENERATED FILE. Edit canonical lessons, then run npm run generate:books.
% canonical-source-hash: fnv1a64:fbb0d944d9c040d7
% canonical-lessons: SA-C136-prapnoti, SA-C136-dharayati, SA-C136-ganayati, SA-C136-marjayati, SA-C136-badhnati

\chapter{Reaches, Wears, Counts, Wipes, Ties}
\label{ch:sa-reaches-wears-counts-wipes-ties}
\hlchaptermodality{136}

\begin{chapteropening}
\textbf{By the end of this chapter:} \emph{I can say reaches, wears, counts, wipes and ties.}

Complete the last lesson of chapter 136: I can say reaches, wears, counts, wipes and ties.
\end{chapteropening}

\section[\texorpdfstring{\sk{प्राप्नोति} (prāpnoti)}{प्राप्नोति (prāpnoti)}]{\sk{प्राप्नोति} (prāpnoti) — reaches, obtains}

\label{lesson:SA-C136-prapnoti}



\begin{quote}
Before the new one: say the Sanskrit for climbs, then the Sanskrit for climbs down.
\end{quote}

\subsection*{You'll want to know: \sk{प्राप्नोति}}
\textbf{\sk{प्राप्नोति}} — \emph{prāpnoti} — \textquotedblleft{}reaches, obtains\textquotedblright{}.

\begin{grammarlens}[title={What you've built}]
One more word for this chapter.
\end{grammarlens}

\subsection*{Your turn}
\begin{itemize}
\raggedright
  \item \emph{Say it:} \emph{prāpnoti}
  \item \emph{Say it:} \emph{prāpnoti}, once more
  \item \emph{Say it:} say \emph{prāpnoti}
  \item \emph{Recall:} say the Sanskrit for climbs, then the Sanskrit for climbs down, then say \emph{prāpnoti} again
\end{itemize}

\begin{tcolorbox}[breakable,colback=teal!4,colframe=teal!35!black,title={Before you move on}]
What does \sk{प्राप्नोति} mean? (\textquotedblleft{}Reaches, obtains\textquotedblright{}.) Say it once more.
\end{tcolorbox}
\section[\texorpdfstring{\sk{धारयति} (dhārayati)}{धारयति (dhārayati)}]{\sk{धारयति} (dhārayati) — wears, holds}

\label{lesson:SA-C136-dharayati}



\begin{quote}
Before the new one: say the Sanskrit for climbs down, then the Sanskrit for reaches.
\end{quote}

\subsection*{You'll want to know: \sk{धारयति}}
\textbf{\sk{धारयति}} — \emph{dhārayati} — \textquotedblleft{}wears, holds\textquotedblright{}.

\begin{grammarlens}[title={What you've built}]
One more word for this chapter.
\end{grammarlens}

\subsection*{Your turn}
\begin{itemize}
\raggedright
  \item \emph{Say it:} \emph{dhārayati}
  \item \emph{Say it:} \emph{dhārayati}, once more
  \item \emph{Say it:} say \emph{dhārayati}
  \item \emph{Recall:} say the Sanskrit for climbs down, then the Sanskrit for reaches, then say \emph{dhārayati} again
\end{itemize}

\begin{tcolorbox}[breakable,colback=teal!4,colframe=teal!35!black,title={Before you move on}]
What does \sk{धारयति} mean? (\textquotedblleft{}Wears, holds\textquotedblright{}.) Say it once more.
\end{tcolorbox}
\section[\texorpdfstring{\sk{गणयति} (gaṇayati)}{गणयति (gaṇayati)}]{\sk{गणयति} (gaṇayati) — counts}

\label{lesson:SA-C136-ganayati}



\begin{quote}
Before the new one: say the Sanskrit for reaches, then the Sanskrit for wears.
\end{quote}

\subsection*{You'll want to know: \sk{गणयति}}
\textbf{\sk{गणयति}} — \emph{gaṇayati} — \textquotedblleft{}counts\textquotedblright{}.

\begin{grammarlens}[title={What you've built}]
One more word for this chapter.
\end{grammarlens}

\subsection*{Your turn}
\begin{itemize}
\raggedright
  \item \emph{Say it:} \emph{gaṇayati}
  \item \emph{Say it:} \emph{gaṇayati}, once more
  \item \emph{Say it:} say \emph{gaṇayati}
  \item \emph{Recall:} say the Sanskrit for reaches, then the Sanskrit for wears, then say \emph{gaṇayati} again
\end{itemize}

\begin{tcolorbox}[breakable,colback=teal!4,colframe=teal!35!black,title={Before you move on}]
What does \sk{गणयति} mean? (\textquotedblleft{}Counts\textquotedblright{}.) Say it once more.
\end{tcolorbox}
\section[\texorpdfstring{\sk{मार्जयति} (mārjayati)}{मार्जयति (mārjayati)}]{\sk{मार्जयति} (mārjayati) — wipes, cleans}

\label{lesson:SA-C136-marjayati}



\begin{quote}
Before the new one: say the Sanskrit for wears, then the Sanskrit for counts.
\end{quote}

\subsection*{You'll want to know: \sk{मार्जयति}}
\textbf{\sk{मार्जयति}} — \emph{mārjayati} — \textquotedblleft{}wipes, cleans\textquotedblright{}.

\begin{grammarlens}[title={What you've built}]
One more word for this chapter.
\end{grammarlens}

\subsection*{Your turn}
\begin{itemize}
\raggedright
  \item \emph{Say it:} \emph{mārjayati}
  \item \emph{Say it:} \emph{mārjayati}, once more
  \item \emph{Say it:} say \emph{mārjayati}
  \item \emph{Recall:} say the Sanskrit for wears, then the Sanskrit for counts, then say \emph{mārjayati} again
\end{itemize}

\begin{tcolorbox}[breakable,colback=teal!4,colframe=teal!35!black,title={Before you move on}]
What does \sk{मार्जयति} mean? (\textquotedblleft{}Wipes, cleans\textquotedblright{}.) Say it once more.
\end{tcolorbox}
\section[\texorpdfstring{\sk{बध्नाति} (badhnāti)}{बध्नाति (badhnāti)}]{\sk{बध्नाति} (badhnāti) — ties}

\label{lesson:SA-C136-badhnati}



\begin{quote}
Before the new one: say the Sanskrit for counts, then the Sanskrit for wipes.
\end{quote}

\subsection*{You'll want to know: \sk{बध्नाति}}
\textbf{\sk{बध्नाति}} — \emph{badhnāti} — \textquotedblleft{}ties\textquotedblright{}.

\begin{grammarlens}[title={What you've built}]
One more word for this chapter.
\end{grammarlens}

\subsection*{Your turn}
\begin{itemize}
\raggedright
  \item \emph{Say it:} \emph{badhnāti}
  \item \emph{Say it:} \emph{badhnāti}, once more
  \item \emph{Say it:} say \emph{badhnāti}
  \item \emph{Recall:} say the Sanskrit for counts, then the Sanskrit for wipes, then say \emph{badhnāti} again
\end{itemize}

\begin{tcolorbox}[breakable,colback=teal!4,colframe=teal!35!black,title={Before you move on}]
What does \sk{बध्नाति} mean? (\textquotedblleft{}Ties\textquotedblright{}.) Say it once more.
\end{tcolorbox}
